Thank you for your valuable questions.

**A1:** Our dataset contains 7,133 LUAD pathological images (1,519 MIP-positive, 5,614 MIP-negative; imbalance ratio $\approx$ 1:3.7). Each image was annotated by a practicing pathologist, reviewed by a senior pathologist ($>10$ years subspecialty experience), and approved by the institutional data management committee. Upon acceptance, we will release the dataset and source code to support reproducible research.

**A2:** Unlike prior selective classification work in computational pathology that focuses on uncertainty rejection or coverage–accuracy trade-offs, our method targets a clinically distinct goal: controlling false negatives for rare subtypes. We achieve this through three integrated components: (i) cost-sensitive training to reshape the decision boundary, (ii) Neyman–Pearson calibration to provide statistical FNR guarantees, and (iii) an FNR-aware selection rule that directly limits missed high-risk cases in the accepted subset. This contrasts with post-hoc thresholding or overall risk optimization without class-asymmetric guarantees, providing verifiable FNR control with selective abstention beyond conventional uncertainty-aware models.

**A3:** While reserving 100 positive samples for NP calibration was feasible in our study (1,519 MIP-positive cases total), we acknowledge this may be challenging for ultra-rare conditions. In the revised manuscript, we will: (1) quantify the minimum positive sample size required for reliable calibration using the binomial bound in Theorem 2.1 (e.g., for $\alpha=5\\%$ FNR control with $\delta=0.05$, $m_1 \geq 59$); and (2) discuss using generative models to synthesize positive samples for calibration, while noting its dependence on sample quality. These discussions will be added to the conclusion.


**A4:** In the revised conclusion, we will add actionable guidelines for selecting the target FNR based on clinical context: for high-risk settings, we recommend a strict 1–3\% FNR target despite lower coverage (40–60\%); for routine screening, a 5\% target offers practical balance (~91\% coverage, deferring ~9 cases per 100 patients, Case 4, $w_1=10$); for resource-constrained workflows, a 10\% target may be acceptable. These recommendations will be directly linked to Tables 6–8 and estimated review workloads. 